\documentclass[10pt,a4paper]{amsart}
\usepackage{amsmath,amssymb,amsthm,graphicx}
% Retained minimal preamble for the Stacks Project chapter flat.tex (More on Flatness).
% Verbatim excerpt of the chapter preamble: Xy-pic setup (native preamble lines 14-18) and the native
% theorem declarations with the mathematical macros (native preamble lines 154-217). Copied mechanically;
% nothing added, removed or redefined. The wrapper class and the amsmath/amssymb/amsthm/graphicx base
% packages are supplied by the assembler, so the class/externaldocument part of the native preamble is omitted.
% For commutative diagrams we use Xy-pic
\usepackage[all]{xy}

% We use 2cell for 2-commutative diagrams.
\xyoption{2cell}
% Theorem environments.
%
\theoremstyle{plain}
\newtheorem{theorem}[subsection]{Theorem}
\newtheorem{proposition}[subsection]{Proposition}
\newtheorem{lemma}[subsection]{Lemma}

\theoremstyle{definition}
\newtheorem{definition}[subsection]{Definition}
\newtheorem{example}[subsection]{Example}
\newtheorem{exercise}[subsection]{Exercise}
\newtheorem{situation}[subsection]{Situation}

\theoremstyle{remark}
\newtheorem{remark}[subsection]{Remark}
\newtheorem{remarks}[subsection]{Remarks}

\numberwithin{equation}{subsection}

% Macros
%
\def\lim{\mathop{\mathrm{lim}}\nolimits}
\def\colim{\mathop{\mathrm{colim}}\nolimits}
\def\Spec{\mathop{\mathrm{Spec}}}
\def\Hom{\mathop{\mathrm{Hom}}\nolimits}
\def\Ext{\mathop{\mathrm{Ext}}\nolimits}
\def\SheafHom{\mathop{\mathcal{H}\!\mathit{om}}\nolimits}
\def\SheafExt{\mathop{\mathcal{E}\!\mathit{xt}}\nolimits}
\def\Sch{\mathit{Sch}}
\def\Mor{\mathop{\mathrm{Mor}}\nolimits}
\def\Ob{\mathop{\mathrm{Ob}}\nolimits}
\def\Sh{\mathop{\mathit{Sh}}\nolimits}
\def\NL{\mathop{N\!L}\nolimits}
\def\CH{\mathop{\mathrm{CH}}\nolimits}
\def\proetale{{pro\text{-}\acute{e}tale}}
\def\etale{{\acute{e}tale}}
\def\QCoh{\mathit{QCoh}}
\def\Ker{\mathop{\mathrm{Ker}}}
\def\Im{\mathop{\mathrm{Im}}}
\def\Coker{\mathop{\mathrm{Coker}}}
\def\Coim{\mathop{\mathrm{Coim}}}

% Boxtimes
%
\DeclareMathSymbol{\boxtimes}{\mathbin}{AMSa}{"02}

%
% Macros for moduli stacks/spaces
%
\def\QCohstack{\mathcal{QC}\!\mathit{oh}}
\def\Cohstack{\mathcal{C}\!\mathit{oh}}
\def\Spacesstack{\mathcal{S}\!\mathit{paces}}
\def\Quotfunctor{\mathrm{Quot}}
\def\Hilbfunctor{\mathrm{Hilb}}
\def\Curvesstack{\mathcal{C}\!\mathit{urves}}
\def\Polarizedstack{\mathcal{P}\!\mathit{olarized}}
\def\Complexesstack{\mathcal{C}\!\mathit{omplexes}}
% \Pic is the operator that assigns to X its picard group, usage \Pic(X)
% \Picardstack_{X/B} denotes the Picard stack of X over B
% \Picardfunctor_{X/B} denotes the Picard functor of X over B
\def\Pic{\mathop{\mathrm{Pic}}\nolimits}
\def\Picardstack{\mathcal{P}\!\mathit{ic}}
\def\Picardfunctor{\mathrm{Pic}}
\def\Deformationcategory{\mathcal{D}\!\mathit{ef}}

% --- appended 2026-10-12 (Chief): the project own definitions of the reference/slogan/history
% display environments, copied mechanically from the shared Stacks preamble lines 8-12
% (\usepackage{verbatim} plus the three \newenvironment lines, which the chapter preamble keeps
% verbatim from the project preamble). Nothing added, removed or redefined: these environments
% carry no mathematical content and the project itself typesets them as comments. Without these
% three lines a counted statement whose own text uses them fails with Environment slogan undefined.
\usepackage{verbatim}
\newenvironment{reference}{\comment}{\endcomment}
\newenvironment{slogan}{\comment}{\endcomment}
\newenvironment{history}{\comment}{\endcomment}
\usepackage{cleveref}
\crefname{originalreference}{original reference}{original references}
\Crefname{originalreference}{Original reference}{Original references}
\makeatletter
\newlabel{situation-mod-injective}{{0AT0}{source}{}{sourcepointer-1}{}}
\newlabel{situation-mod-injective@cref}{{[originalreference][0][]0AT0}{[1][0][]source}{}{}{}}
\makeatother
\crefname{originallemma}{original lemma}{original lemmas}
\Crefname{originallemma}{Original lemma}{Original lemmas}
\makeatletter
\newlabel{lemma-reduce-finite-type-injective-into-flat-mod-m}{{0AT2}{source}{}{sourcepointer-2}{}}
\newlabel{lemma-reduce-finite-type-injective-into-flat-mod-m@cref}{{[originallemma][0][]0AT2}{[1][0][]source}{}{}{}}
\makeatother
\crefname{originallemma}{original lemma}{original lemmas}
\Crefname{originallemma}{Original lemma}{Original lemmas}
\makeatletter
\newlabel{lemma-flat-finite-type-local-colimit-free}{{05U7}{source}{}{sourcepointer-3}{}}
\newlabel{lemma-flat-finite-type-local-colimit-free@cref}{{[originallemma][0][]05U7}{[1][0][]source}{}{}{}}
\makeatother
\crefname{originallemma}{original lemma}{original lemmas}
\Crefname{originallemma}{Original lemma}{Original lemmas}
\makeatletter
\newlabel{lemma-universally-injective-local}{{05FQ}{source}{}{sourcepointer-4}{}}
\newlabel{lemma-universally-injective-local@cref}{{[originallemma][0][]05FQ}{[1][0][]source}{}{}{}}
\makeatother
\crefname{originaltheorem}{original theorem}{original theorems}
\Crefname{originaltheorem}{Original theorem}{Original theorems}
\makeatletter
\newlabel{theorem-flattening-local}{{05PG}{source}{}{sourcepointer-5}{}}
\newlabel{theorem-flattening-local@cref}{{[originaltheorem][0][]05PG}{[1][0][]source}{}{}{}}
\makeatother
\crefname{originalreference}{original reference}{original references}
\Crefname{originalreference}{Original reference}{Original references}
\makeatletter
\newlabel{algebra-remark-Tor-ring-mod-ideal}{{00M6}{source}{}{sourcepointer-6}{}}
\newlabel{algebra-remark-Tor-ring-mod-ideal@cref}{{[originalreference][0][]00M6}{[1][0][]source}{}{}{}}
\makeatother
\crefname{originallemma}{original lemma}{original lemmas}
\Crefname{originallemma}{Original lemma}{Original lemmas}
\makeatletter
\newlabel{algebra-lemma-what-does-it-mean}{{051C}{source}{}{sourcepointer-7}{}}
\newlabel{algebra-lemma-what-does-it-mean@cref}{{[originallemma][0][]051C}{[1][0][]source}{}{}{}}
\makeatother
\title{The injectivity lemma holds over a henselian local ring}
\author{The Stacks Project Authors}
\date{stacks-project:flat.tex:local-snapshot-20261011; GNU FDL 1.2 or later}
\begin{document}
\maketitle
\noindent\textit{Editorial note.} All delivered mathematical text is the author's own, copied from the pinned Stacks Project chapter file flat.tex; no statement or proof step was rewritten, reordered or removed.\par\medskip
\noindent\textit{Editorial note.} The preamble is the chapter's own retained minimal preamble (sources/stacks-flat/retained-preamble.tex), a verbatim excerpt of the native Xy-pic setup and of the native theorem declarations and mathematical macros, followed by the project own definitions of the reference, slogan and history display environments (verbatim package plus three newenvironment lines) that the retained file carries at its end and that the project itself typesets as comments. No macro was invented or redefined, and the wrapper class and the amsmath/amssymb/amsthm/graphicx base packages are supplied by the assembler.\par\medskip
\noindent\textit{Editorial note.} The counted statement and proof contain no cross reference, no citation and no included file, so the delivered body is self contained.\par\medskip
\noindent\textit{Editorial note.} Canonical location of the counted proposition: Stacks tag 0AT3 (sources/stacks-flat/pinned-original-source.tex line 7073).\par\medskip
\noindent\textit{Allowed standard prerequisites (editorial).}\begin{itemize}
\item Reduction to the essentially finite presentation case. Same chapter, Lemma labelled lemma-reduce-finite-type-injective-into-flat-mod-m, cited at proof lines 7078-7080. Allows the flattening machinery to be applied to B and N.
\item Flat modules as filtered colimits of free modules with injective transition maps modulo the maximal ideal, and universal injectivity of the colimit. Same chapter, Lemma labelled lemma-flat-finite-type-local-colimit-free, cited at proof lines 7084-7086, and Lemma labelled lemma-universally-injective-local, cited at proof lines 7088-7090. Settles the case where N is flat over A.
\item Local flattening: the ideal measuring the failure of flatness. Same chapter, Theorem labelled theorem-flattening-local, cited at proof lines 7096-7100. Produces the ideal I with N/IN flat and N/I\textasciicircum{}2N not flat unless I = 0.
\item Tor vanishing and flatness criteria. Algebra, Remark labelled algebra-remark-Tor-ring-mod-ideal and Lemma labelled algebra-lemma-what-does-it-mean, cited at proof lines 7118-7122. Converts the injectivity of the lower left arrow into flatness of N/I\textasciicircum{}2N over A/I\textasciicircum{}2.
\end{itemize}\medskip
\ifdefined\hypertarget\hypertarget{sourcepointer-1}{}\fi
\noindent\textit{Source reference.} Original reference 0AT0: sources/stacks-flat/pinned-original-source.tex, source lines 7017--7017. Same-chapter pointer into the author environment situation at source line 7017; invoked at proof line None of the counted proof. It is retained only as a navigation pointer; every load-bearing step is inside the delivered span.\par\medskip
\ifdefined\hypertarget\hypertarget{sourcepointer-2}{}\fi
\noindent\textit{Source reference.} Original lemma 0AT2: sources/stacks-flat/pinned-original-source.tex, source lines 7045--7045. Same-chapter pointer into the author environment lemma at source line 7045; invoked at proof line 7081 of the counted proof. It is retained only as a navigation pointer; every load-bearing step is inside the delivered span.\par\medskip
\ifdefined\hypertarget\hypertarget{sourcepointer-3}{}\fi
\noindent\textit{Source reference.} Original lemma 05U7: sources/stacks-flat/pinned-original-source.tex, source lines 5759--5759. Same-chapter pointer into the author environment lemma at source line 5759; invoked at proof line 7086 of the counted proof. It is retained only as a navigation pointer; every load-bearing step is inside the delivered span.\par\medskip
\ifdefined\hypertarget\hypertarget{sourcepointer-4}{}\fi
\noindent\textit{Source reference.} Original lemma 05FQ: sources/stacks-flat/pinned-original-source.tex, source lines 1946--1946. Same-chapter pointer into the author environment lemma at source line 1946; invoked at proof line 7089 of the counted proof. It is retained only as a navigation pointer; every load-bearing step is inside the delivered span.\par\medskip
\ifdefined\hypertarget\hypertarget{sourcepointer-5}{}\fi
\noindent\textit{Source reference.} Original theorem 05PG: sources/stacks-flat/pinned-original-source.tex, source lines 6790--6790. Same-chapter pointer into the author environment theorem at source line 6790; invoked at proof line 7095 of the counted proof. It is retained only as a navigation pointer; every load-bearing step is inside the delivered span.\par\medskip
\ifdefined\hypertarget\hypertarget{sourcepointer-6}{}\fi
\noindent\textit{Source reference.} Original reference 00M6: sources/stacks-chow/tags.txt, source lines 812--812. Cross-chapter standard prerequisite located in the project tag list staged in the local snapshot; invoked at proof line 7118 of the counted proof. It is retained only as a navigation pointer; every load-bearing step is inside the delivered span.\par\medskip
\ifdefined\hypertarget\hypertarget{sourcepointer-7}{}\fi
\noindent\textit{Source reference.} Original lemma 051C: sources/stacks-chow/tags.txt, source lines 6530--6530. Cross-chapter standard prerequisite located in the project tag list staged in the local snapshot; invoked at proof line 7120 of the counted proof. It is retained only as a navigation pointer; every load-bearing step is inside the delivered span.\par\medskip
\begin{lemma}
\label{lemma-henselian-finite-type-injective-into-flat-mod-m}
If in Situation \ref{situation-mod-injective} the ring $A$ is henselian
then the lemma holds.
\end{lemma}

\begin{proof}
It suffices to prove this when $B$ is essentially of finite presentation
over $A$ and $N$ is of finite presentation over $B$, see
Lemma \ref{lemma-reduce-finite-type-injective-into-flat-mod-m}.
Let us temporarily make the additional assumption that $N$ is flat over $A$.
Then $N$ is a filtered colimit $N = \colim_i F_i$
of free $A$-modules $F_i$ such that the transition maps
$u_{ii'} : F_i \to F_{i'}$ are injective modulo $\mathfrak m_A$, see
Lemma \ref{lemma-flat-finite-type-local-colimit-free}.
Each of the compositions $u_i : F_i \to M$ is $A$-universally
injective by
Lemma \ref{lemma-universally-injective-local}
wherefore $u = \colim u_i$ is $A$-universally injective as desired.

\medskip\noindent
Assume $A$ is a henselian local ring, $B$ is essentially
of finite presentation over $A$, $N$ of finite presentation over $B$. By
Theorem \ref{theorem-flattening-local}
there exists a finitely generated ideal $I \subset A$ such that
$N/IN$ is flat over $A/I$ and such that $N/I^2N$ is not flat over
$A/I^2$ unless $I = 0$. The result of the previous paragraph shows that
the lemma holds for $u \bmod I : N/IN \to M/IM$ over $A/I$.
Consider the commutative diagram
$$
\xymatrix{
0 \ar[r] &
M \otimes_A I/I^2 \ar[r] &
M/I^2M \ar[r] &
M/IM \ar[r] & 0 \\
&
N \otimes_A I/I^2 \ar[r] \ar[u]^u &
N/I^2N \ar[r] \ar[u]^u &
N/IN \ar[r] \ar[u]^u & 0
}
$$
whose rows are exact by right exactness of $\otimes$ and the fact that
$M$ is flat over $A$. Note that the left vertical arrow is the map
$N/IN \otimes_{A/I} I/I^2 \to M/IM \otimes_{A/I} I/I^2$, hence is
injective. A diagram chase shows that the lower left arrow is injective,
i.e., $\text{Tor}^1_{A/I^2}(I/I^2, M/I^2) = 0$ see
Algebra, Remark \ref{algebra-remark-Tor-ring-mod-ideal}.
Hence $N/I^2N$ is flat over $A/I^2$ by
Algebra, Lemma \ref{algebra-lemma-what-does-it-mean}
a contradiction unless $I = 0$.
\end{proof}

\end{document}
